% EXTRACTED from v1.1 manuscript (tag v1.1-verified): Part X, ch. "Modal,
% Temporal, and Epistemic Logic" sec. "Normal propositional modal logic"
% (definition of the Hilbert calculus K), cross-checked against the "Rule
% Sheets" chapter (sec:rule-sheets-03).  Metavariables $A,B$ follow the
% nonclassical-part convention.
\calculussheet{sheet:hilbert-k}{Modal Hilbert Calculus $K$}{Language: modal propositional formulas --- the classical propositional grammar ($\bot,\neg,\land,\lor,\to$ primitive) extended by the constructor $\Box A$, with $\Diamond A\defeq\neg\Box\neg A$ defined; sequent shape: no sequents --- Hilbert-style derivability $\Gamma\derives_{K}A$ via finite sequences whose lines are members of $\Gamma$, axiom-schema instances, modus-ponens consequences of earlier lines, or formulas $\Box A$ for which $A$ has a separate no-premise derivation (necessitation is theorem-only even in a derivation from premises); negation: primitive (classical); equality: none.}
\hbrule{rule:hilbert-k-taut}{Taut}{}{A}{$A$ is a uniform substitution instance of a classically valid propositional formula}
\hbrule{rule:hilbert-k-k}{K}{}{\Box(A\to B)\to(\Box A\to\Box B)}{\none}
\hbrule{rule:hilbert-k-mp}{MP}{A\to B\qquad A}{B}{\none}
\hbrule{rule:hilbert-k-nec}{Nec}{\derives_{K}A}{\derives_{K}\Box A}{theorem-only: applies when $A$ has a derivation from no premises, not to an arbitrary formula depending on an undischarged premise}
